%! Author = lili
%! Date = 7/15/26

% Zellzyklus, Kontrolle der Zellteilung
% 36 Folien
% Chapter 9 (parts)
%… to the Cell Cycle

\section*{Mitosestadien}
\defin{prophase}
\defin{prometaphase}
\defin{metaphase}
\defin{anaphase}
\defin{telophase}
\defin{cytokinese}

% TODO reihenfolge der phasen
% TODO phasen  der mitose aufzählen
\section{Der Zellzyklus}
\defin{interphase}
Sie Interphase ist unterteilt in mehrere Abschnitte:
% TODO phasen der interphase aufzählen
% was passiert in G1
% was passiert in S-Phase
% Was passiert in G2
% für jede Phase: wie ist die ploidie? 2n? 4n? etc.

Nach Abschluß einer Mitose beginnt die G1-Phase.
Es findet die RNA- und Proteinsynthese statt, aber keine DNA-Replikation.

Der Beginn der S-Phase ist
markiert durch Initiation der
DNA-Replikation.
Sie endet mit der
vollständigen Replikation der
gesamten DNA.

Die G2-Phase geht vom Ende der S-Phase bis zum
Beginn der Mitose.

\subsection{Anregung von Zellen zur Zellteilung - Wachstumsfaktor Signal-Transduktionsweg}
\defin{wachstumsfaktor}
Die Bindung des Liganden
``Wachstumsfaktor'' (GF) an den
GF-Rezeptor führt zur
Dimerisierung des Rezeptors und
nachfolgend zur Phosphorylierung
der cytoplasmatischen Domäne
des Rezeptors.
\defin{exchange_factor}
Der GF-Rezeptor bindet den
``\gls{exchange_factor}'' (SOS als GEF)
und aktiviert das G-Protein RAS.
Bindung über das Adapter-
Protein Grb2.

ras ist in Tumoren oft mutiert
(oft gain-of-function Mutationen,
mutiertes Protein liegt permanent
als aktives RAS-GTP vor).

\defin{ras}

% TODO F 15

\subsection{Kontrolle der Progression des Zellzyklus}
% TODO F 16 - 23



%%%%%%%%%%%%%%%%%%%%%%%%%%%%%%%%%%%%%%%%%%%%%%%%%%%%%%%%%%%%%%

\section{Zellzyklus}\label{sec:zellzyklus}

\textbf{Zellzyklus (somatisch):} Zeitraum vom Ende einer Mitose bis zum Ende der nächsten.

\begin{itemize}
    \item \textbf{\gls{interphase}}: Zeitraum zwischen zwei Mitosen (G1 -- S -- G2)
    \item \textbf{M-Phase}: eigentliche Zellteilung; Mitosestadien lichtmikroskopisch sichtbar
    \item Typische Dauern (Säugerzelle):
    \begin{itemize}
        \item G1: 6--12 h (2n)
        \item S: 6--8 h ($2n \rightarrow 4n$)
        \item G2: 3--4 h (4n)
        \item M: ca.
        1 h
    \end{itemize}
    \item \textbf{G0}: Austritt aus dem Zellzyklus (reversibel: \emph{reactivation} oder unbefristet: \emph{indefinite withdrawal})
\end{itemize}
\defin{interphase}

% ABBILDUNG:
% Kreisdiagramm des Zellzyklus mit G1/S/G2/M sowie G0-Abzweig,
% einschließlich DNA-Gehalt (2n/4n) und Zeitangaben.

\subsection{Checkpoints (Kontrollpunkte)}\label{subsec:checkpoints-(kontrollpunkte)}

Kontrollschleifen, bei denen die Initiation eines Abschnitts von der erfolgreichen
Beendigung des vorherigen Abschnitts abhängig ist.

\begin{itemize}
    \item Checkpoint for division (G1)
    \item Checkpoints for DNA integrity (S)
    \item Checkpoints for unreplicated DNA (G2)
    \item Checkpoint for paired kinetochores (M)
\end{itemize}

% ABBILDUNG:
% Schema ``basic cycle control'' vs. ``checkpoint controls''
% entlang G1--S--G2--M.

\subsection{RB1 im Zellzyklus}\label{subsec:rb1-im-zellzyklus}
\defin{cyclin}
\begin{itemize}
    \item \textbf{Restriction Point}: Punkt in G1, an dem die Entscheidung zur Zellteilung fällt.
    \item Das Protein \textbf{Rb} integriert Informationen über DNA-Schäden und Zellwachstum.
    \item Rb bindet die \textbf{E2F}-Transkriptionsfaktoren, wodurch deren Zielgene inaktiv bleiben.
    \item \gls{cyclin}/CDK-Komplexe
    \begin{itemize}
        \item \gls{cyclin} D1/CDK4
        \item \gls{cyclin} D1/CDK6
        \item \gls{cyclin} E/CDK2
    \end{itemize}
    phosphorylieren Rb. Dadurch wird E2F freigesetzt und der Zellzyklus kann fortgesetzt werden.
\end{itemize}

\subsection{Retinoblastom}\label{subsec:retinoblastom}

Retinoblastom ist ein bösartiger Tumor der Netzhaut, der aus unreifen Netzhautzellen hervorgeht.

\begin{itemize}
    \item Beide \gls{allel}e des \textit{RB1}-\gls{gen}s müssen mutiert sein $\rightarrow$ auf Zellebene rezessiv.
    \item Die erbliche, bilaterale Form erscheint phänotypisch dominant, da Träger heterozygot sind und das zweite \gls{wildtyp}-
    \gls{allel} häufig somatisch mutiert.
    \item Dieses Prinzip wird als \textbf{Zwei-Mutationen-Theorie} nach Knudson bezeichnet.
    \item Das gleiche Prinzip gilt für die sporadische, unilaterale Form, bei der beide Mutationen somatisch entstehen.
\end{itemize}

% ABBILDUNG:
% ``RB is a recessive disease'':
% Darstellung der Genotypen RB+/RB+, RB+/RB- und RB-/RB-
% mit entsprechender Tumorbildung.


